\documentclass[12pt]{article}
\usepackage[margin=1in]{geometry}
\usepackage{enumitem}
%\usepackage{forest}
\usepackage{dirtree}
\usepackage{fontawesome5}
\usepackage[utf8]{inputenc}
\usepackage{multicol}
\usepackage[T1]{fontenc}
\usepackage{tabularx}
\usepackage{fancyhdr}
\usepackage{titlesec}
\usepackage{lmodern}
\usepackage{helvet}
\usepackage[breaklinks=true]{hyperref}
\usepackage[svgnames]{xcolor}
\usepackage[table]{xcolor}
\usepackage[framemethod=tikz]{mdframed}
\usepackage[customcolors]{hf-tikz}
\usepackage[most]{tcolorbox}
\usetikzlibrary{shadows}
\usepackage{tgheros}
\renewcommand{\familydefault}{\sfdefault}
\usepackage{listings}
\usepackage[spanish]{babel} 
\usepackage{xcolor,graphicx}
\usepackage{tikz}
\usetikzlibrary{arrows,shapes,positioning,shadows,trees}
%%%%%%%%%%%%%%%%%%%%%%%%%%%%%%%%%%%%%%%%%
%colors
\definecolor{rositapastel}{RGB}{255, 135, 159}
\titleformat{\section}
  {\color{rositapastel}\normalfont\Large\bfseries}
  {\thesection}{1em}{}

\titleformat{\subsection}
  {\color{rositapastel}\normalfont\large\bfseries}
  {\thesubsection}{1em}{}

\titleformat{\subsubsection}
  {\color{rositapastel}\normalfont\normalsize\bfseries}
  {\thesubsubsection}{1em}{}
%%%%%%%%%%%%%%%%%%%%%%
% Mis comandos
\newcounter{recordatorioCounter}
\newcommand{\recordatorio}{%
  \stepcounter{recordatorioCounter}% avanza el contador
  \textbf{\textcolor{rositapastel}{Recordatorio \therecordatorioCounter : }}%
}



\lstdefinestyle{mystyle}{
    backgroundcolor=\color{rositapastel!23},   
    commentstyle=\color{codegreen},
    keywordstyle=\color{magenta},
    numberstyle=\tiny\color{codegray},
    stringstyle=\color{codepurple},
    basicstyle=\ttfamily\footnotesize,
    breakatwhitespace=false,         
    breaklines=true,                 
    captionpos=b,                    
    keepspaces=true,                 
    numbers=left,                   
    numbersep=5pt,                  
    showspaces=false,                
    showstringspaces=false,
    showtabs=false,                  
    tabsize=2
}

\lstset{style=mystyle}

% Icono Carpeta
\newcommand{\foldericon}{\faFolderOpen~}

\definecolor{foldercolor}{RGB}{255, 193, 7}
\definecolor{pdfcolor}{RGB}{220, 53, 69}
\definecolor{readmecolor}{RGB}{40, 167, 69}

\renewcommand{\foldericon}{\textcolor{foldercolor}{\faFolder}}
\newcommand{\pdficon}{\textcolor{pdfcolor}{\faFilePdf}}
\newcommand{\txticon}{\textcolor{txtcolor}{\faFileAlt}}
\newcommand{\zipicon}{\textcolor{gray}{\faFileArchive}}



% Set up fancy header/footer
\pagestyle{fancy}
\fancyhead[LO,L]{UNAM}
\fancyhead[RO,R]{\textit{\textsc{Proyecto}}}   
\fancyfoot[LO,L]{Facultad De Ciencias, 2026-I}
\fancyfoot[CO,C]{\thepage}
\fancyfoot[RO,R]{Modelado y Programación}
\renewcommand{\headrulewidth}{0.4pt}
\renewcommand{\footrulewidth}{0.4pt}
%%%%%%%%%%%%%%%%%%%%%%

\newdimen\digitwidth
\settowidth\digitwidth{0}
\def~{\hspace{\digitwidth}}

% Cambiar color de secciones
\titleformat{\section}
  {\color{rositapastel}\normalfont\Large\bfseries}
  {\thesection}{1em}{}

\titleformat{\subsection}
  {\color{rositapastel}\normalfont\large\bfseries}
  {\thesubsection}{1em}{}

\titleformat{\subsubsection}
  {\color{rositapastel}\normalfont\normalsize\bfseries}
  {\thesubsubsection}{1em}{}

% Cambiar color de ítems en listas numeradas
\setlist[enumerate]{label=\textcolor{rositapastel}{\arabic*.}, left=1.5em}
\setlist[itemize]{label=\textcolor{rositapastel}{\alph*}, left=1.5em}
\setlist[itemize]{label=\textcolor{rositapastel}{\textbullet}, left=1.5em}


\begin{document}
\begin{titlepage}
\begin{center}

\textsc{\Large\uppercase{Universidad Nacional Autónoma De México}}\\[0.4cm]	

{\huge \Large \uppercase{Facultad de Ciencias, 2027-I} \\[0.4cm] }
{\huge \Large \uppercase{Facultad de Ciencias} \\[1cm] }

%%\includegraphics[scale=.2]{image.png}\\[0.6cm]


%%% Título
\rule{\linewidth}{0.4pt}\\[0.7em]
{ \color{rositapastel} \textsc{\Large Proyecto:}} \\[0.7em] 
 \textit{\Large Chat}
 \\[0.1em] 
\rule{\linewidth}{0.4pt}\\[2em]


\color{black!80!black}{\large \uppercase{\textsc{ Profesor:}}}\\[0.25cm]

\begin{tabular}{l}
  \large Canek Pelaez Valdez \\[0.4cm]
\end{tabular}

\large \uppercase{Ayudante de Teoría:}\\[0.2cm]
\centering
\begin{tabular}{l}

\large Daniel Rojo Mata \\[0.4cm]
\end{tabular}

\large \uppercase{Ayudante de Laboratorio:}\\[0.2cm]
\centering
\begin{tabular}{l}

\large Irvin Javier Cruz Gonzalez\\[0.4cm]
\end{tabular}

\vfill

\end{center}
\end{titlepage}

\section*{Resumen}
La idea del Proyecto 1 es crear un chat en línea en el que múltiples clientes se puedan conectar y mandar mensajes entre ellos, con otras diversas operaciones dentro de el chat, la idea es tener una sala general para todos los usuarios y que además cada uno pueda hacer sus salas privadas, invitando a otros usuarios, también podrán solicitar la lista de usuarios conectados, tanto para la sala general como para las que son parte, actualizar su estado, y desconectarse del chat.
En este reporte vamos a ver primero una introducción con una descripción del proyecto, luego vamos a ver el proceso de analisis y modelado, hasta llegar a la implementación y posteriormente explicar las conclusiones.

\section*{Introduccion}
Este proyecto es







\newpage
\section*{Analisis y diseño}

E chat se divide en dos partes, cliente y servidor:

\subsection*{Cliente}

El cliente es el que representa al usuario y debe tener ciertas operaciones como:



  
\newpage
        \begin{itemize}
\item
          \begin{center}
          \begin{tikzpicture}[
            every node/.style = {
              draw,
              circle,
              minimum size=1.5cm
            },
            level 1/.style = {sibling distance=4cm},
            edge from parent/.style = {draw, -latex}
          ]
          \node[draw,,rounded corners] {Raiz}
            child {
              node[fill=red!20] {}
              edge from parent node[left,draw=none,pos=0.1] {0}
            }
            child {
              node[fill=rositapastel!70] {A}
              edge from parent node[right,draw=none,pos=0.1] {1}
            };
          \end{tikzpicture}
          \end{center}   
          
        \textbf{RECUERDA} en las ramas de la izquierda va un cero y en las de la derecha va 1 y se coloca la letra con mayor frecuencia de repetición.

        \item A partir de aquí esto será un proceso iterativo, donde siempre tomaremos el nodo izquierdo para derivar 2 ramas más y colocar en la rama de la derecha la siguiente letra con mayor frecuencia, por lo tanto nuestra segunda aproximación del árbol binario es esta:


            \begin{center}
            \begin{tikzpicture}[
              every node/.style = {
                draw,
                circle,
                minimum size=1.5cm
              },
              level 1/.style = {sibling distance=4cm},
              level 2/.style = {sibling distance=3cm},
              edge from parent/.style = {draw, -latex}
            ]
            \node[draw,,rounded corners] {Raiz}
              child {
                node [fill=red!20] {}
                child {
                  node[fill=red!20] {}
                  edge from parent node[left,draw=none,pos=0.1] {0}
                }
                child {
                  node[fill=rositapastel!70] {R}
                  edge from parent node[right,draw=none,pos=0.1] {1}
                }
                edge from parent node[left,draw=none,pos=0.1] {0}
              }
              child {
                node[fill=rositapastel!70] {A}
                edge from parent node[right,draw=none,pos=0.1] {1}
              };
            \end{tikzpicture}
            \end{center}

 \newpage       
  %%     Podemos observar que las demás letras su frecuencia es la misma en este caso no importada si colocamos primero una y despúes otra. Repetimos los mismos pasos con el resto de las letras y finalmente tenemos el árbol final resultante:
      
  %%     \begin{center}
  %%     \begin{tikzpicture}[scale=0.8,
  %%       every node/.style = {
  %%         draw,
  %%         circle,
  %%         minimum size=1cm,
  %%         scale=0.8,
  %%         transform shape
  %%       },
  %%       level 1/.style = {sibling distance=6cm},
  %%       level 2/.style = {sibling distance=4cm},
  %%       level 3/.style = {sibling distance=3cm},
  %%       edge from parent/.style = {draw, -latex}
  %%     ]
  %%     \node[draw,,rounded corners] {Raiz}
  %%       child {
  %%         node[fill=red!20] {}
  %%         child {
  %%           node[fill=red!20] {}
  %%           child {
  %%             node[fill=red!20] {}
  %%             child {
  %%               node[fill=red!20] {}
  %%               child {
  %%                 node[fill=red!20] {}
  %%                 child {
  %%                   node[fill=red!20] {}
  %%                   child {
  %%                     node[fill=red!20] {}
  %%                     edge from parent node[left,draw=none,pos=0.3] {0}
  %%                   }
  %%                   child {
  %%                     node[fill=antiquefuchsia!70] {B}
  %%                     edge from parent node[right,draw=none,pos=0.3] {1}
  %%                   }
  %%                   edge from parent node[left,draw=none,pos=0.3] {0}
  %%                 }
  %%                 child {
  %%                   node[fill=antiquefuchsia!70] {D}
  %%                   edge from parent node[right,draw=none,pos=0.3] {1}
  %%                 }
  %%                 edge from parent node[left,draw=none,pos=0.3] {0}
  %%               }
  %%               child {
  %%                 node[fill=antiquefuchsia!70] {K}
  %%                 edge from parent node[right,draw=none,pos=0.3] {1}
  %%               }
  %%               edge from parent node[left,draw=none,pos=0.3] {0}
  %%             }
  %%             child {
  %%               node[fill=antiquefuchsia!70] {H}
  %%               edge from parent node[right,draw=none,pos=0.3] {1}
  %%             }
  %%             edge from parent node[left,draw=none,pos=0.3] {0}
  %%           }
  %%           child {
  %%             node[fill=antiquefuchsia!70] {V}
  %%             edge from parent node[right,draw=none,pos=0.3] {1}
  %%           }
  %%           edge from parent node[left,draw=none,pos=0.3] {0}
  %%         }
  %%         child {
  %%           node[fill=antiquefuchsia!70] {R}
  %%           edge from parent node[right,draw=none,pos=0.3] {1}
  %%         }
  %%         edge from parent node[left,draw=none,pos=0.3] {0}
  %%       }
  %%       child {
  %%         node[fill=antiquefuchsia!70] {A}
  %%         edge from parent node[right,draw=none,pos=0.3] {1}
  %%       };
  %%     \end{tikzpicture}
  %%     \end{center}

   \end{itemize}
 
  
  \subsubsection*{  Recorriendo el árbol.}
   
  Ahora que ya tenemos construido nuestro árbol, podemos recorrerlo para buscar una letra, por ejemplo, si queremos encontrar a la letra 'V' tendríamos que hacer el siguiente recorrido marcado en color rojo:

      % No marca el recorrido a una sola hoja sino a todo el árbol. Tu tarea será encontrar la solución.

      %%      \begin{center}
      %% \begin{tikzpicture}[scale=0.8,
      %%   every node/.style = {
      %%     draw,
      %%     circle,
      %%     minimum size=1cm,
      %%     scale=0.8,
      %%     transform shape
      %%   },
      %%   level 1/.style = {sibling distance=6cm},
      %%   level 2/.style = {sibling distance=4cm},
      %%   level 3/.style = {sibling distance=3cm},
      %%   edge from parent/.style = {draw, -latex},
      %%   highlightpath/.style = {edge from parent/.style = {draw=red, very thick, -latex}}
      %% ]
      %% \node[draw,,rounded corners] {Raiz}
      %%   child[highlightpath] {
      %%     node[fill=red!20] {}
      %%     child[highlightpath] {
      %%       node[fill=red!20] {}
      %%       child[missing]
      %%       child[highlightpath] {
      %%         node[fill=red!20] {V}
      %%         edge from parent node[right,draw=none,pos=0.3] {1}
      %%       }
      %%       edge from parent node[left,draw=none,pos=0.3] {0}
      %%     }
      %%     child {
      %%       node[fill=antiquefuchsia!70] {R}
      %%       edge from parent node[right,draw=none,pos=0.3] {1}
      %%     }
      %%     edge from parent node[left,draw=none,pos=0.3] {0}
      %%   }
      %%   child {
      %%     node[fill=antiquefuchsia!70] {A}
      %%     edge from parent node[right,draw=none,pos=0.3] {1}
      %%   };
      %% \end{tikzpicture}
      %% \end{center}

  
  
      Como podemos observar, vemos representado por 3 bits y no por 8, para llegar a la letra 'V' siguiendo la ruta marcada en rojo que corresponde a la secuencia '001', es decir, 
      primero tomamos la rama izquierda (0), luego otra rama izquierda (0), y finalmente la rama derecha (1).

      Para 'A' solo se hace un simple recorrido y se obtiene 1 como se representación binaria, que igual reducimos de 8 a 1 bit.        


      \begin{table}[h]
            \centering
            \begin{tabular}{|c|c|c|}
                \hline
                Letra & Frecuencia & Representación\\
                \hline
                 A & 5 & 1 \\
                \hline
                 R & 2 & 01 \\
                \hline
                 V & 1 & 001 \\
                \hline
                 H & 1 & 0001 \\
                \hline
                 K & 1 & 00001\\
                \hline
                 D & 1  & 000001 \\
                \hline
                 B & 1 & 0000001\\
                \hline
            \end{tabular}
            \caption{Representación binaria del resto de letras}
            \label{tab:frecuencias}
        \end{table}

   \subsubsection*{Nueva representación binaria} 
   
    \begin{center}
      \texttt{1001011000100001100000110000001011}
    \end{center}    
       
      \begin{center}
%%          \includegraphics[scale=.35]{abra.png}
        \end{center}

        Hemos pasado de 96 bits a 34

  \subsection*{Huffman Decoding} 

  El decodificador debe saber qué hay de inicio un archivo comprimido, leerlo y construir el árbol de Huffman para el alfabeto. Solo entonces puede leer y decodificar el resto de su entrada. Es decir para este algoritmo debe recibir una \textbf{cadena de bits} y \textbf{el árbol de Huffman de esa cadena}. \\


  El algoritmo comienza:
  \begin{enumerate}
    \item Desde la raiz del árbol, lee el primer bit de entrada (del texto ya comprimido), si es cero recorre la rama izquierda del árbol, si es 1 recorre la rama derecha.
    \item Recorre el árbol de acuerdo al bit hasta llegar a un nodo hoja.
    \item Ya que llegó a un hoja genera el símbolo original. Repite hasta procesar todos los bits.
  \end{enumerate} 

\newpage
%% \subsubsection*{Ejemplo - Decoding}

%% Recibimos el arbol ya construido y la representación binaria de la cadena ya codificada/comprimida 

%%    \begin{multicols}{2}
%%       \begin{center}
%%       \begin{tikzpicture}[scale=0.55,
%%         every node/.style = {
%%           draw,
%%           circle,
%%           minimum size=1cm,
%%           scale=0.8,
%%           transform shape
%%         },
%%         level 1/.style = {sibling distance=6cm},
%%         level 2/.style = {sibling distance=4cm},
%%         level 3/.style = {sibling distance=3cm},
%%         edge from parent/.style = {draw, -latex}
%%       ]
%%       \node[draw,,rounded corners] {Raiz}
%%         child {
%%           node[fill=red!20] {}
%%           child {
%%             node[fill=red!20] {}
%%             child {
%%               node[fill=red!20] {}
%%               child {
%%                 node[fill=red!20] {}
%%                 child {
%%                   node[fill=red!20] {}
%%                   child {
%%                     node[fill=red!20] {}
%%                     child {
%%                       node[fill=red!20] {}
%%                       edge from parent node[left,draw=none,pos=0.3] {0}
%%                     }
%%                     child {
%%                       node[fill=antiquefuchsia!70] {B}
%%                       edge from parent node[right,draw=none,pos=0.3] {1}
%%                     }
%%                     edge from parent node[left,draw=none,pos=0.3] {0}
%%                   }
%%                   child {
%%                     node[fill=antiquefuchsia!70] {D}
%%                     edge from parent node[right,draw=none,pos=0.3] {1}
%%                   }
%%                   edge from parent node[left,draw=none,pos=0.3] {0}
%%                 }
%%                 child {
%%                   node[fill=antiquefuchsia!70] {K}
%%                   edge from parent node[right,draw=none,pos=0.3] {1}
%%                 }
%%                 edge from parent node[left,draw=none,pos=0.3] {0}
%%               }
%%               child {
%%                 node[fill=antiquefuchsia!70] {H}
%%                 edge from parent node[right,draw=none,pos=0.3] {1}
%%               }
%%               edge from parent node[left,draw=none,pos=0.3] {0}
%%             }
%%             child {
%%               node[fill=antiquefuchsia!70] {V}
%%               edge from parent node[right,draw=none,pos=0.3] {1}
%%             }
%%             edge from parent node[left,draw=none,pos=0.3] {0}
%%           }
%%           child {
%%             node[fill=antiquefuchsia!70] {R}
%%             edge from parent node[right,draw=none,pos=0.3] {1}
%%           }
%%           edge from parent node[left,draw=none,pos=0.3] {0}
%%         }
%%         child {
%%           node[fill=antiquefuchsia!70] {A}
%%           edge from parent node[right,draw=none,pos=0.3] {1}
%%         };
%%       \end{tikzpicture}
%%       \end{center}

%%       \begin{center}
%%       \texttt{1001011000100001100000110000001011}
%%       \end{center}
%%    \end{multicols}


   \begin{enumerate}
    \item Tomamos el primer bit de la cadena \texttt{(1)} y recorremos el arbol desde la raiz, notamos que \texttt{1} es la rama derecha que tiene como hoja la letra 'A'. 
          Entonces la agregamos a una lista.

          \texttt{\textcolor{red}{1}001011000100001100000110000001011}

         \begin{center}
    %%      \includegraphics[scale=.14]{a.jpg}
        \end{center}

          \begin{center}
            [A]
          \end{center}
           
\newpage
    \item Pasamos al segundo bit que es \texttt{(0)} y hacemos el recorrido al arbol desde la raiz, vemos que no es hoja, asi que nos tomamos el siguiente bit \texttt{(0)} que tampoco es hoja, nos volvemos a tomar el siguiente bit \texttt{(1)} que es hoja y contiene a 'V'.
          Entonces 'V' la agregamos a la lista que teniamos 

        \texttt{1\textcolor{red}{001}011000100001100000110000001011}  

           \begin{center}
  %%        \includegraphics[scale=.14]{v.jpg}
        \end{center}

          \begin{center}
            [AV] 
          \end{center}

    \item Pasamos al siguiente bit despues de los bits que nos genera la anteriores letras. que en este caso es \texttt{01}. Recorremos el árbol y vemos que el cero no es hoja, seguimos recorriendo y vemos que 1 ya es hoja, así que agregamos a la lista al elemento que contiene, 'R'.
    
      \texttt{1001\textcolor{red}{01}1000100001100000110000001011}

     \begin{center}
      %%    \includegraphics[scale=.14]{r.jpg}
        \end{center}
    
    \begin{center}
            [AVR] 
          \end{center} 
          
   \end{enumerate}

   El proceso sigue hasta procesar todos los bits, agregando sus respectivas letras a la lista:
   \begin{center}
    \texttt{[AVRAHKADABRA]}
   \end{center} 
      


  


\section*{Actividades - Documentación}
\begin{enumerate}
  \item ¿Cuál es la diferencia entre la compresión con perdida y compresión sin perdida?
  \item Investiga como se lleva a cabo la compresión de texto, images, videos y audio.
  \item ¿A que se refiere el enfoque para la compresión de datos, códigos de longitud variable?
  \item Crea 2 archivos de texto donde:
        \begin{enumerate}
          \item [a)] Ingresa una tus frases favoritas SIN ESPACIOS y revisa su contenido binario mediante el comando \texttt{xxd} y su peso en bytes.
          \item [b)] Ingresa  la misma frase ahora CON ESPACIOS y revisa su contenido binario mediante el comando \texttt{xxd} y su peso en bytes.
        \end{enumerate}
        ¿Cuál es la diferencia respecto a contenido binario de los dos archivos?

  \item Comprime la frase del anterior inciso y otra frase más, usando los algoritmos vistos. Para el primer algoritmo adjunta su representación binaria de inicio y final despúes de utilizarlo.  
  \item ¿Por qué es necesario tener un respectivo árbol de Huffman para decodificar una cadena de texto? ¿Qué pasa si no lo tengo?     
  \item ¿De qué manera las estrategias de optimización empleadas en el videojuego \textit{Grand Theft Auto V} para consolas de septima generación (PS3/Xbox 360) reflejan los principios fundamentales de la compresión de datos en comparación con las de octava generación?      
\end{enumerate}


\section*{Actividades - Implementación}
\begin{enumerate}
  \item Dada una cadena de texto, regresar la frecuencias de cada uno de sus caracteres.
  \item Proporciona el arbol de Huffman asociado a una cadena de texto.
  \item Implementa el algoritmo para codificar y comprimir una cadena de texto.
  \item Implementa el algoritmo para decoficar un cadena de texto ya comprimida.
  \item Agrega un ejemplo de cada algoritmo.
\end{enumerate}


\section*{Reporte}

\begin{itemize}
  \item Reporte detallado del proyecto (pdf) donde se considere:
        \begin{itemize}
          \item \textbf{Portada}. Nombre y número de cuenta del alumno
          \item \textbf{Introducción.} Descripción del proyecto, objetivo y la importancia de conocer esta información.
          \item \textbf{Desarrollo.} Explicar bajo que estrategias o enfoques se llevó a cabo el proyecto documentado paso a paso, además de cuales son las entradas y salidas estandar. Tambien en este criterio se llevan a cabo las actividades
          \item \textbf{Conclusiones.} Como se pueden mejorar estos algoritmos, da otro enfoque (puede ser breve) desde la perspectiva de la programación imperativa, comentarios de lo que se llego a aprender, cuales fueron las dificultades y como se resolvieron al momento de implementar.
          \item \textbf{Referencias.} Fuentes en formato \textit{APA 7}, se prohibe el uso de Wikipedia.
        \end{itemize}
\end{itemize}


\section*{Material de Consulta}

\begin{itemize}

    \item Video: Huffman Coding Explained – Computerphile
          \begin{center}
          \textcolor{rositapastel}{\url{https://www.youtube.com/watch?v=JsTptu56GM8}}
          \end{center}

     \item De texto a Binario
          \begin{center}
          \textcolor{rositapastel}{\url{https://www.rapidtables.com/convert/number/ascii-to-binary.html}}
          \end{center}

     \item Implementación con JavaScript y Java.
       \begin{center}
          \textcolor{rositapastel}{\url{https://github.com/martacanirome4/HuffmanCoding?tab=readme-ov-file}}
        \end{center}       
       
  \end{itemize}          



\section*{Lineamientos para proyecto }
\begin{itemize}
    \item Se deben realizar \href{https://medium.com/@hfigueroa.dev/commits-at%C3%B3micos-f7e02434130}{\textcolor{red}{commits átomicos}}. Despúes de completar una funcionalidad o caracteristica realizar un commit, esto se debe de hacer periodicamente a lo largo del tiempo establecido del proyecto. (evaluable)
    \item El proyecto debe contar con un archivo \texttt{README.md} donde se indique: Nombre del proyecto, objetivo, funcionamiento, ejecución, el tiempo requerido, y si es necesario comentarios extra.  
     \item Toda función debe estar comentada.
     \item Se debe hacer uso de modulos.
     \item Prohibido el uso de \texttt{fold}.
    \item Se debe de hacer uso de commits semánticos.
    \item  Prohibido el uso de cualquier IA, la entrega de esta proyecto será sometida a la prueba Turing.
    
\end{itemize}

\section*{Entregables.}

Deberás compartir lo siguiente mediante  \textit{Google Classroom}:

\begin{itemize}
    \item Enlace a tu repositorio de \textit{GitHub}.
    \item Reporte en formato pdf
    \item Captura de pantalla de tu terminal al ejecutar el comando \texttt{git log} despúes de terminar tu proyecto. 
\end{itemize}

Tu repositorio debe estar organizado de la siguiente manera:\\


\dirtree{%
.1 \foldericon\ EstructurasDiscretas.
.2 \foldericon\ Practica1/.
.3 \faIcon[regular]{file-alt}\ README.md.
.3  \faIcon[regular]{file-alt}\ practica1.md.
.2 \foldericon\ Practica2/.
.3 \faIcon[regular]{file-alt}\ practica2.hs.
.3 \faIcon[regular]{file-alt}\ README.md.
.2 \foldericon\ Practica3/.
.3 \faIcon[regular]{file-alt}\ practica3.hs.
.3 \faIcon[regular]{file-alt}\ README.md.
.2 \foldericon\ Practica4/.
.3 \faIcon[regular]{file-alt}\ natural.hs.
.3 \faIcon[regular]{file-alt}\ mati.hs.
.3 \faIcon[regular]{file-alt}\ Aux.hs.
.3 \faIcon[regular]{file-alt}\ README.md.
.2 \foldericon\ Practica5/.
.3 \faIcon[regular]{file-alt}\ Practica5.hs.
.3 \faIcon[regular]{file-alt}\ Aux.hs.
.3 \faIcon[regular]{file-alt}\ README.md.
.2 \foldericon\ Practica6/.
.3 \faIcon[regular]{file-alt}\ Practica6.hs.
.3 \faIcon[regular]{file-alt}\ Aux.hs.
.3 \faIcon[regular]{file-alt}\ README.md.
.2 \foldericon\ Proyecto/.
.3 \faIcon[regular]{file-alt}\ \textcolor{rositapastel!110}{Proyecto.hs}.
.3 \faIcon[regular]{file-alt}\ \textcolor{rositapastel!110}{Aux.hs}.
.3 \faIcon[regular]{file-alt}\ \textcolor{rositapastel!110}{README.md}.
.2  \faIcon[regular]{file-alt}\ README.md.
.2  \faIcon[regular]{file-alt}\ LICENSE.
}








\end{document}


